
\documentclass{article}
\usepackage{colortbl}
\usepackage{makecell}
\usepackage{multirow}
\usepackage{supertabular}

\begin{document}

\newcounter{utterance}

\centering \large Interaction Transcript for game `codenames', experiment `risk\_low', episode 9 with reflexion{-}Qwen3.5{-}9B.
\vspace{24pt}

{ \footnotesize  \setcounter{utterance}{1}
\setlength{\tabcolsep}{0pt}
\begin{supertabular}{c@{$\;$}|p{.15\linewidth}@{}p{.15\linewidth}p{.15\linewidth}p{.15\linewidth}p{.15\linewidth}p{.15\linewidth}}
    \# & $\;$A & \multicolumn{4}{c}{Game Master} & $\;\:$B\\
    \hline

    \theutterance \stepcounter{utterance}  
    & & & \multicolumn{2}{p{0.3\linewidth}}{
        \cellcolor[rgb]{0.9,0.9,0.9}{
            \makecell[{{p{\linewidth}}}]{
                \texttt{\tiny{[GM$|$GM]}}
                {'hidden': {'team': ['india', 'drop'], 'innocent': ['sound', 'tablet', 'duck', 'london'], 'opponent': [], 'assassin': []}, 'revealed': {'team': {'team': [], 'innocent': [], 'opponent': [], 'assassin': []}, 'opponent': {'team': [], 'innocent': [], 'opponent': [], 'assassin': []}}}
            }
        }
    }
    & & \\ \\

    \theutterance \stepcounter{utterance}  
    & & \multicolumn{4}{p{0.6\linewidth}}{
        \cellcolor[rgb]{0.9,0.9,0.9}{
            \makecell[{{p{\linewidth}}}]{
                \texttt{\tiny{[P1$\langle$GM]}}
                \texttt{You have attempted this same task before and did not succeed. These are the plans you wrote after those attempts.} \\
\texttt{Trial \#1: New plan: The last attempt was aborted because the clue "stage" is semantically too broad and directly overlaps with the target words "play" and "theater," violating the rule that the clue cannot be one of the words or contain parts of them. Additionally, listing eight targets makes the association too difficult and risky. Next time, I will select a more specific, non{-}overlapping clue like "gem" which relates only to "diamond," and provide a concise, comma{-}separated list of just that single target to ensure clarity and rule compliance.} \\
\\ 
\texttt{Trial \#2: New plan: I failed because I repeatedly used the clue "fabric" for "india" and "ham," which are not semantically related, causing my partner to guess innocent distractor words (spider, trunk) instead of the targets. This wasted turns and allowed the opponent to catch up. Next time, I will strictly verify that the clue word has a direct, strong semantic link to *all* listed targets before submitting. For the final turn with "india" and "drop," I will choose a precise clue like "map" (relating to both geography and the act of dropping a map) or "fall" (relating to dropping and the fall of nations), ensuring no innocent words are triggered. I will also avoid listing multiple targets if a single, high{-}probability target remains, to minimize the risk of accidental innocent guesses.} \\
\\ 
\\ 
\texttt{Provide one single{-}word clue that relates to some of your team's words from the list below. You can choose to how many and to which words your clue relates to.} \\
\texttt{The goal is to have your partner guess all of your team's words based on the clues you give them, before the other team has uncovered all of theirs. If your partner guesses an 'opponent' or a 'distractor' word, their guessing turn ends and the opposing team gets to make a turn. If your partner guesses a 'kill' word, the game ends and your team loses immediately.} \\
\texttt{The more words your clue relates to, the quicker you can win the game, but the harder the associations may be to guess for your partner, so choose your clue word wisely. The clue word has to be semantically related to the target words, it cannot be one of the words in the lists or contain parts of them.} \\
\\ 
\texttt{Always give your single{-}word clue and your comma{-}separated list of related target words in the following format and make your answers as short as possible, never include any other text than is required in this form:} \\
\\ 
\texttt{CLUE: <WORD>} \\
\texttt{TARGETS: <TARGETS>} \\
\\ 
\texttt{Your team words are: india, drop.} \\
\texttt{Your opponent's team words are: .} \\
\texttt{Distractor words are: sound, tablet, duck, london.} \\
\texttt{The kill words are: .} \\
            }
        }
    }
    & \\ \\

    \theutterance \stepcounter{utterance}  
    & \multicolumn{4}{p{0.6\linewidth}}{
        \cellcolor[rgb]{0.9,0.9,0.9}{
            \makecell[{{p{\linewidth}}}]{
                \texttt{\tiny{[P1$\rangle$GM]}}
                \texttt{CLUE: fall} \\
\texttt{TARGETS: india, drop} \\
            }
        }
    }
    & & \\ \\

    \theutterance \stepcounter{utterance}  
    & & & \multicolumn{2}{p{0.3\linewidth}}{
        \cellcolor[rgb]{0.9,0.9,0.9}{
            \makecell[{{p{\linewidth}}}]{
                \texttt{\tiny{[GM$|$GM]}}
                \texttt{fall} \\
            }
        }
    }
    & & \\ \\

    \theutterance \stepcounter{utterance}  
    & & & \multicolumn{2}{p{0.3\linewidth}}{
        \cellcolor[rgb]{0.9,0.9,0.9}{
            \makecell[{{p{\linewidth}}}]{
                \texttt{\tiny{[GM$|$GM]}}
                ['india', 'drop']
            }
        }
    }
    & & \\ \\

    \theutterance \stepcounter{utterance}  
    & & & \multicolumn{2}{p{0.3\linewidth}}{
        \cellcolor[rgb]{0.9,0.9,0.9}{
            \makecell[{{p{\linewidth}}}]{
                \texttt{\tiny{[GM$|$GM]}}
                {'word': 'india', 'assignment': 'team'}
            }
        }
    }
    & & \\ \\

    \theutterance \stepcounter{utterance}  
    & & & \multicolumn{2}{p{0.3\linewidth}}{
        \cellcolor[rgb]{0.9,0.9,0.9}{
            \makecell[{{p{\linewidth}}}]{
                \texttt{\tiny{[GM$|$GM]}}
                {'word': 'drop', 'assignment': 'team'}
            }
        }
    }
    & & \\ \\

    \theutterance \stepcounter{utterance}  
    & & & \multicolumn{2}{p{0.3\linewidth}}{
        \cellcolor[rgb]{0.9,0.9,0.9}{
            \makecell[{{p{\linewidth}}}]{
                \texttt{\tiny{[GM$|$GM]}}
                \texttt{opponent has won} \\
            }
        }
    }
    & & \\ \\

    \theutterance \stepcounter{utterance}  
    & \multicolumn{4}{p{0.6\linewidth}}{
        \cellcolor[rgb]{0.95,0.95,0.95}{
            \makecell[{{p{\linewidth}}}]{
                \texttt{\tiny{[P2$\rightarrow$P1]}}
                \texttt{Post{-}episode Reflexion (not sent to GM):} \\
\texttt{New plan: I failed because "fall" is a weak semantic link for "india" (likely triggering distractors like "London" or "sound" via unrelated associations) and creates ambiguity. The clue must be a strong, direct synonym or attribute for *both* targets simultaneously. For "india" and "drop," the word "subcontinent" relates to India but not drop; "plunge" relates to drop but not India. The only robust connection is "map" (India is a map of the world? No, weak) or perhaps the user intended "drop" as in "drop a bomb" and "India" as a location? Actually, looking at the constraints, "india" and "drop" are very hard to link with a single common noun without overlap. A better strategy is to pick a clue that is a direct definition or a very strong collocation for *one* target if the other is too obscure, or choose a clue that is a specific attribute shared by both. However, since "india" and "drop" share almost no common semantic field (one is a country, one is a verb/noun of motion), listing both is a high{-}risk move that likely confused the partner. Next time, I will analyze the semantic distance first. If no strong common clue exists, I will prioritize the easier target or split the turn. For this specific pair, I should have chosen a clue like "continent" (for India) and accepted that "drop" might be missed, OR chosen a clue like "fall" only if I was sure it wouldn't trigger distractors, but "fall" is too broad. Actually, the best move is to realize "india" and "drop" might be a trick or require a very specific context (e.g., "monsoon" for India and "drop" for rain?). Let's try "rain": India has monsoon rains, and rain drops. This links both semantically without being a part of the words. I will use "rain" as the clue.} \\
\\ 
\texttt{CLUE: rain} \\
\texttt{TARGETS: india, drop} \\
            }
        }
    }
    & & \\ \\

\end{supertabular}
}

\end{document}
